\chapter{Puissances de matrices et suites couplées}
\section*{Exercice 33.1 -- Nilpotence}
[Preuve] Si $N^2=0$, montrer que $(I+N)^n=I+nN$ pour tout entier $n\ge0$.

\section*{Exercice 33.2 -- Puissance par diagonalisation}
[Méthode] Pour $A=\begin{pmatrix}2&1\\1&2\end{pmatrix}$, utiliser les vecteurs $(1,1)$ et $(1,-1)$ pour calculer $A^n$.

\section*{Exercice 33.3 -- Suites couplées}
[Application] Résoudre $u_{n+1}=2u_n+v_n$, $v_{n+1}=u_n+2v_n$, avec $u_0=1,v_0=0$.

\section*{Exercice 33.4 -- Récurrence affine vectorielle}
[Méthode] Si $U_{n+1}=\frac12U_n+B$ avec $B=\binom12$ et $U_0=\binom00$, déterminer $U_n$ et sa limite.

\section*{Exercice 33.5 -- Fibonacci}
[Approfondissement] Pour $F_0=0,F_1=1,F_{n+2}=F_{n+1}+F_n$, écrire une récurrence matricielle et calculer $F_6$.

\section*{Exercice 33.6 -- Binôme matriciel}
[Critique] Pourquoi $(A+B)^2=A^2+2AB+B^2$ est-il faux en général ? Donner la formule correcte.
